\documentclass[10pt,a4paper]{amsart}
\usepackage[margin=19mm]{geometry}
\usepackage{amsmath,amssymb,mathtools}
\usepackage[scr=rsfs,cal=euler]{mathalfa}
\usepackage{xfrac}
\usepackage[unicode,hidelinks,bookmarks=false]{hyperref}
\newcommand{\ie}{i.e.\ }
\newcommand{\cf}{cf.\ }
\newcommand{\resp}{respectively\ }
\newcommand{\Subsection}{\subsection}
\providecommand{\nobreakdash}{\nobreak\hbox{-}}
\setlength{\emergencystretch}{2em}
\multlinegap0pt
\usepackage{bm}
\usepackage{bbm}
\usepackage{dsfont}
\usepackage{xspace}
\usepackage{cancel}
\usepackage{mathtools}
\usepackage{amsthm}
\makeatletter
\@ifundefined{theorem}{\newtheorem{theorem}{Theorem}}{}
\@ifundefined{proposition}{\newtheorem{proposition}[theorem]{Proposition}}{}
\makeatother
\title[On SAGBI bases theory of skew Poincar\'e-Birkhoff-Witt exten]{On SAGBI bases theory of skew Poincar\'e-Birkhoff-Witt extensions}
\author{Y\'esica Su\'arez, Armando Reyes}
\date{Source: arXiv:2508.10715v1 (CC BY 4.0)}
\begin{document}
\maketitle

\noindent\emph{Source and licence.} On SAGBI bases theory of skew Poincar\'e-Birkhoff-Witt extensions. Version arXiv:2508.10715v1 (CC BY 4.0), distributed under the Creative Commons Attribution 4.0 International licence, as recorded in the arXiv licence metadata for this exact version and shown on the abstract page https://arxiv.org/abs/2508.10715v1. Full rights notice: licenses/arxiv-2508.10715-CC-BY-4.0.txt.\par\smallskip

\noindent\textit{Editorial scope.} Counted unit: the Theorem whose source label is \texttt{sagbicriterio} in the article (source0.tex lines 853--855, source label \texttt{sagbicriterio}), delivered together with the article's own complete original proof of it (source0.tex lines 856--889, one \texttt{proof} environment of 34 lines). No mathematics is written, repaired or re-derived: statement and proof are the article's own text line for line. Required context retained uncounted: prop (lines 805--812); 1 (lines 807--811). Every definition, display and label of those lines is retained unchanged. Omitted from this package and not redistributed: 1--804, 812--852, 890--1249. Nothing inside a retained excerpt is omitted or altered: every retained body is delivered exactly as the article's lines are written, so no omission receipt of any kind accompanies this package. Citations: the retained excerpts carry no citation command at all, so no bibliography entry of the article is transcribed and no reference is redistributed. Reference adaptation: none was needed and none was made; every reference inside the delivered unit resolves inside it, so no unresolved reference or citation is produced. Printed slips kept literal: no printed word, symbol, hypothesis or number of the counted statement or of its proof was altered. Layout and class: the article's own class is \texttt{amsart} with packages including amsthm, amssymb, amsfonts, latexsym, mathtools, thmtools, mathrsfs, fontenc, tikz-cd, enumitem, hyperref, algorithm2e, xy; the wrapper is the lane's pinned \texttt{amsart} editorial wrapper at 10pt on A4, which restates the theorem-like environments the retained text needs and adds the article's own macro definitions that the retained text needs (none), with extra wrapper packages (bm, bbm, dsfont, xspace, cancel, mathtools). The delivered numbering is the unit's own. Assets: no retained span contains a figure environment, a graphics-inclusion command, a PDF-page inclusion, a table or a TikZ picture; no image file is copied into this package, the package declares no asset dependency and no image file is redistributed with it. Licence: version arXiv:2508.10715v1 of this article is distributed under the Creative Commons Attribution 4.0 International licence, as recorded in the arXiv licence metadata for this exact version and shown on the abstract page https://arxiv.org/abs/2508.10715v1. A full rights notice is delivered at licenses/arxiv-2508.10715-CC-BY-4.0.txt. Characters: every retained excerpt is pure ASCII, so the delivered engine, XeLaTeX with its default Latin Modern text font, reports no missing character.
\begin{proposition}\label{prop}
Let $S$ be subalgebra of a SPBW extension $A$ and $F \subseteq S$. If $F$ is a {\rm \texttt{SAGBI}} basis of $S$, then the following assertions hold:
\begin{enumerate}
\item[\rm{(1)}] \label{1} For each  $s\in A$, we have that $s \in S$ if and only {\rm \texttt{SNF}}$(s | F) = 0$.

\item[\rm{(2)}]  $F$ generates the subalgebra $S$, i.e. $S=\Bbbk\langle F\rangle_{A}$.
\end{enumerate}
\end{proposition}

\begin{theorem}\label{sagbicriterio} 
Suppose that $F$ generates $S$ as a subalgebra in the SPBW extension $A.$ Then $F$ is a {\rm {\texttt{SAGBI}}} basis of $S$ if an only if every $T$-polynomial of every critical pair of $F$ gives zero {\rm {\texttt{SAGBI}}} normal form.
\end{theorem}

\begin{proof}
If $H$ is a \texttt{SAGBI} basis of $S$ then every $T$-polynomial is an element of $S=\Bbbk\langle F\rangle_{A}$ and its \texttt{SAGBI} normal form is equal to zero by part \ref{1} of Proposition \ref{prop}.

Conversely, suppose given $0 \neq s \in S .$ It is sufficient to prove that it has a representation 
$$
s = \sum\limits_{p=1}^{t} k_{p} m_{p}(F),
$$ 

where $k_{p} \in \Bbbk$ and $m_{p}(F) \in \Bbbk\langle F\rangle_{A}$ with 
$$
\operatorname{lm}(s)=\operatorname{ht}\left(\sum\limits_{p=1}^{t} k_{p} m_{p}(F)\right)=\operatorname{lm}(m_i(F)).
$$

Let $s \in S$ given by $s = \sum\limits_{p=1}^{t} k_{p} m_{p}(F)$ with smallest possible height $X$ among all possible representations of $s$ in $S,$ that is, $X=\max _{p=1}^{t}\left\{\operatorname{lm}\left(m_{p}(F)\right)\right\}$. It is clear that $ X \succeq\operatorname{lm}(s)$.

Suppose that $X\succ\operatorname{lm}(s)$ i.e. cancellation of terms occur then there exist at least two $F$-monomials such that their leading monomial is equal to $X .$ Assume we have only two $F$-monomials $m_{i}(F), m_{j}(F)$ in the representation 
$$
s = \sum\limits_{p=1}^{t} k_{p} m_{p}(F)
$$ 

such that 
$$
\operatorname{lm}\left(m_{i}(F)\right)=\operatorname{lm}\left(m_{j}(F)\right)=X.
$$ 

If $T\left(m_{i}(F), m_{j}(F)\right)=m_{i}(F)-k m_{j}(F)$, then 
\begin{align}
s &=\sum_{p=1}^{t} k_{p} m_{p}(F) \\
&=k_{i}\left(m_{i}(F)-k m_{j}(F)\right) +\left(k_{j}+k_{i} k\right) m_{j}(F) \ +\sum_{p=1, p \neq i, j}^{t} k_{p} m_{p}(F)\\
&=k_{i} T\left(m_{i}(F), m_{j}(F)\right)+\left(k_{j}+k_{i} k\right) m_{j}(F) \ +\sum_{p=1, p \neq i, j}^{t} k_{p} m_{p}(F).\label{ec10}
\end{align}

Since $T\left(m_{i}(F), m_{j}(F)\right)$ has a zero \texttt{SAGBI} normal form, this $T$-polynomial is either zero or can be written as sum of $F$-monomials of height $\operatorname{lm}\left(T\left(m_{i}(F), m_{j}(F)\right)\right)$ which is less than $X .$ If $k_{j}+k_{i} k$ is equal to zero, then the right-hand side of expression (\ref{ec10}) is a representation of $s$ with height less than $X,$ which contradicts our initial assumption that we have chosen a representation of $s$ with smallest possible height. Otherwise, the height is preserved, but on the right-hand side of expression (\ref{ec10}), we have only one $F$-monomial $m_{j}(F)$ such that $\operatorname{lm}\left(m_{j}(F)\right)=X,$ which is a contradiction as at least two $F$-monomials of such type must exist in the representation of $s$.
\end{proof}

\end{document}
